\begin{table}[htbp]
\centering
\caption{Where the computed library loses, resolved by spectral region. The measured-to-computed cosine is compared with the measured-to-measured cosine in the same region, which is the ceiling that measurement repeatability permits. The last four columns split the same comparison by whether the agent contains phosphorus (seven agents) or not (three).}
\label{tab:si-regions}
\footnotesize
\resizebox{\textwidth}{!}{%
\begin{tabular}{lrrrrrrrr}
\toprule
Region (\wn) & Computed & Repeatability & Gap & Intensity & P-containing & No P & $\Delta$ & $p$ \\
\midrule
250--600 & 0.770 & 0.905 & 0.136 & 20.2\% & 0.810 & 0.675 & -0.135 & 0.967 \\
\textbf{600--850} & \textbf{0.719} & \textbf{0.982} & \textbf{0.263} & \textbf{29.2\%} & \textbf{0.636} & \textbf{0.912} & \textbf{+0.276} & \textbf{0.058} \\
850--1150 & 0.826 & 0.982 & 0.156 & 16.6\% & 0.840 & 0.793 & -0.047 & 0.808 \\
1150--1400 & 0.796 & 0.982 & 0.186 & 12.4\% & 0.814 & 0.755 & -0.058 & 0.867 \\
1400--1800 & 0.698 & 0.989 & 0.291 & 21.6\% & 0.758 & 0.557 & -0.201 & 1.000 \\
\bottomrule
\end{tabular}
}

\begin{flushleft}\footnotesize The largest gap coincides with the largest share of measured intensity. Within the 600--850\,\wn{} window the deficit belongs to the phosphorus-containing agents; in each of the four other windows the ordering reverses. With three agents on one side and seven on the other, the smallest attainable one-sided $p$ is 0.008 and the observed value is $0.058$, so the separation is not significant at the 0.05 level; the argument rests on the size of the effect and on its confinement to one window rather than on that $p$ value. $p$ is one-sided Mann--Whitney.\end{flushleft}
\end{table}
